\documentclass[11pt]{article}
\usepackage[utf8]{inputenc}
\usepackage{geometry}
\usepackage{booktabs}
\usepackage{amsmath}
\usepackage{hyperref}
\usepackage{enumitem}
\geometry{margin=1in}

\title{Beyond Stochasticism: An Ontological Framework for Agentic Stability}
\author{BitConcepts Research}
\date{\today}

\begin{document}
\maketitle

\begin{abstract}
This paper presents the OEA (Ontology, Epistemic, Agentic) framework as a testable protocol for improving reliability in agentic systems under recursive synthetic-data exposure. Rather than assuming collapse is universally inevitable, we frame OEA as an intervention hypothesis grounded in three layers: ontological constraint definition, epistemic falsification and calibration, and closed-loop empirical action. The manuscript positions OEA relative to model-collapse literature, retrieval-grounded generation work, and calibration methods, then specifies experiments designed to measure whether OEA reduces degradation and improves rejection of synthetic noise.
\end{abstract}

\section{Introduction}
Model collapse has become a central concern in recursive synthetic training pipelines, with evidence of degradation under specific feedback-loop assumptions \cite{shumailov2024modelcollapse,fu2025theoretical}. At the same time, subsequent studies show collapse behavior depends strongly on data-mixing assumptions, curation policy, and accumulation dynamics rather than being a single universal outcome \cite{gerstgrasser2024inevitable,schaeffer2025position}. This motivates a controlled framework that separates what is known, what is hypothesized, and what is empirically testable.

The OEA framework is proposed as such a structure. \textbf{Ontology} defines the boundary conditions of valid outputs. \textbf{Epistemic} introduces explicit falsification and calibration checks for generated claims. \textbf{Agentic} operationalizes interventions in measurable experiment loops. Philosophically, this aligns with a Popperian criterion: claims must be framed so they can fail under evidence \cite{popper1959logic}.

\section{Related Work and Evidence Positioning}
\subsection{Recursive Synthetic Training and Collapse Dynamics}
Empirical and theoretical analyses report failure modes in recursive loops, including loss of tail behavior and quality drift \cite{shumailov2024modelcollapse,fu2025theoretical}. Newer analyses also identify mitigation regimes, especially when real data are retained or accumulated \cite{gerstgrasser2024inevitable}. Therefore, we treat collapse as regime-dependent rather than monolithic.

\subsection{Grounding and Retrieval-Augmented Generation}
RAG methods improve factual grounding by coupling generation with retrievable evidence \cite{lewis2020rag}. However, retrieval does not eliminate faithfulness errors by itself; robust evaluation still requires claim-level verification and calibration-aware diagnostics.

\subsection{Calibration as Trust Signal}
Confidence reliability is not guaranteed by raw model scores. Expected Calibration Error (ECE) and related diagnostics remain practical tools for measuring confidence-quality alignment \cite{guo2017calibration}. OEA uses calibration as a monitored variable, not a proxy for truth.

\subsection{Trust-Aware Orchestration in Agentic Pipelines}
Recent trust-aware orchestration proposals provide implementation motifs for confidence-mediated arbitration and retrieval-triggered re-evaluation loops \cite{roumeliotis2025trust}. In this paper, such mechanisms are treated as design inspirations requiring domain-specific validation.

\section{Methodology}
\subsection{OEA Tri-Layer Protocol}
\begin{enumerate}[leftmargin=1.5em]
  \item \textbf{Ontological Anchoring}: define non-negotiable domain constraints, variable definitions, and invalid-output criteria before inference.
  \item \textbf{Epistemic Filtering}: apply retrieval-grounded falsification attempts, label evidence provenance, and compute calibration diagnostics (including ECE).
  \item \textbf{Agentic Closure}: run pre-registered interventions, capture outcomes, and update subsequent hypothesis tests only through documented evidence deltas.
\end{enumerate}

\subsection{Primary Experiments}
\textbf{E1: Recursive Stability Test.} Initialize a fixed knowledge seed; run 10 recursive generations where output at step $t-1$ becomes input at step $t$. Compare control vs OEA-conditioned pipelines on divergence and semantic retention metrics.

\textbf{E2: Epistemic Friction Test.} Inject plausible synthetic falsehoods into context and measure rejection/acceptance rates. Report both false accepts and false rejects to avoid one-sided scoring.

\subsection{Pre-Registered Reporting Contract}
For each run we require:
\begin{itemize}
  \item configuration fingerprint (model version, prompts, retrieval settings, seed),
  \item machine-readable artifacts (CSV/JSON),
  \item interval estimates for key metrics,
  \item failure-case exemplars with evidence traces.
\end{itemize}

\section{Results Status and Analysis Plan}
This repository version defines the analysis contract and citation-locked framing; it does not claim completed experimental outcomes yet. Final publication claims are gated on:
\begin{itemize}
  \item executed E1/E2 runs with reproducible artifacts,
  \item confidence intervals and effect-size reporting,
  \item consistency checks between manuscript claims and recorded outputs.
\end{itemize}

\section{Discussion}
The central claim under test is \emph{conditional}: ontological and epistemic constraints may improve stability in specific recursive regimes. We explicitly avoid stronger claims (e.g., universal prevention of collapse) unless supported by cross-regime evidence. This safeguards against over-generalization while preserving actionable engineering hypotheses.

\section{Limitations and Threats to Validity}
\begin{itemize}
  \item \textbf{External validity}: outcomes may not transfer across model families or domains.
  \item \textbf{Retrieval dependence}: evidence quality is bounded by corpus coverage and retrieval precision.
  \item \textbf{Metric insufficiency}: single metrics (including ECE) are non-exhaustive; multiple diagnostics are required.
  \item \textbf{Preprint risk}: some cited work is preprint-stage and may change after peer review.
\end{itemize}

\section{Reproducibility and Ethics Notes}
All major claims should map to reproducible artifacts and cited sources. For deployment-facing use, trust scores and confidence estimates must not be treated as guarantees of truth, especially in safety-critical pathways.

\bibliographystyle{plain}
\bibliography{references}
\end{document}
